\documentclass[10pt,a4paper]{amsart}
\usepackage{amsmath,amssymb,amsthm,graphicx}
% For commutative diagrams we use Xy-pic
\usepackage[all]{xy}

% We use 2cell for 2-commutative diagrams.
\xyoption{2cell}
% Theorem environments.
%
\theoremstyle{plain}
\newtheorem{theorem}[subsection]{Theorem}
\newtheorem{proposition}[subsection]{Proposition}
\newtheorem{lemma}[subsection]{Lemma}

\theoremstyle{definition}
\newtheorem{definition}[subsection]{Definition}
\newtheorem{example}[subsection]{Example}
\newtheorem{exercise}[subsection]{Exercise}
\newtheorem{situation}[subsection]{Situation}

\theoremstyle{remark}
\newtheorem{remark}[subsection]{Remark}
\newtheorem{remarks}[subsection]{Remarks}

\numberwithin{equation}{subsection}

% Macros
%
\def\lim{\mathop{\mathrm{lim}}\nolimits}
\def\colim{\mathop{\mathrm{colim}}\nolimits}
\def\Spec{\mathop{\mathrm{Spec}}}
\def\Hom{\mathop{\mathrm{Hom}}\nolimits}
\def\Ext{\mathop{\mathrm{Ext}}\nolimits}
\def\SheafHom{\mathop{\mathcal{H}\!\mathit{om}}\nolimits}
\def\SheafExt{\mathop{\mathcal{E}\!\mathit{xt}}\nolimits}
\def\Sch{\mathit{Sch}}
\def\Mor{\mathop{\mathrm{Mor}}\nolimits}
\def\Ob{\mathop{\mathrm{Ob}}\nolimits}
\def\Sh{\mathop{\mathit{Sh}}\nolimits}
\def\NL{\mathop{N\!L}\nolimits}
\def\CH{\mathop{\mathrm{CH}}\nolimits}
\def\proetale{{pro\text{-}\acute{e}tale}}
\def\etale{{\acute{e}tale}}
\def\QCoh{\mathit{QCoh}}
\def\Ker{\mathop{\mathrm{Ker}}}
\def\Im{\mathop{\mathrm{Im}}}
\def\Coker{\mathop{\mathrm{Coker}}}
\def\Coim{\mathop{\mathrm{Coim}}}

% Boxtimes
%
\DeclareMathSymbol{\boxtimes}{\mathbin}{AMSa}{"02}

%
% Macros for moduli stacks/spaces
%
\def\QCohstack{\mathcal{QC}\!\mathit{oh}}
\def\Cohstack{\mathcal{C}\!\mathit{oh}}
\def\Spacesstack{\mathcal{S}\!\mathit{paces}}
\def\Quotfunctor{\mathrm{Quot}}
\def\Hilbfunctor{\mathrm{Hilb}}
\def\Curvesstack{\mathcal{C}\!\mathit{urves}}
\def\Polarizedstack{\mathcal{P}\!\mathit{olarized}}
\def\Complexesstack{\mathcal{C}\!\mathit{omplexes}}
% \Pic is the operator that assigns to X its picard group, usage \Pic(X)
% \Picardstack_{X/B} denotes the Picard stack of X over B
% \Picardfunctor_{X/B} denotes the Picard functor of X over B
\def\Pic{\mathop{\mathrm{Pic}}\nolimits}
\def\Picardstack{\mathcal{P}\!\mathit{ic}}
\def\Picardfunctor{\mathrm{Pic}}
\def\Deformationcategory{\mathcal{D}\!\mathit{ef}}
\usepackage{cleveref}
\crefname{originallemma}{original lemma}{original lemmas}
\Crefname{originallemma}{Original lemma}{Original lemmas}
\makeatletter
\newlabel{more-algebra-lemma-henselization}{{0A02}{source}{}{sourcepointer-1}{}}
\newlabel{more-algebra-lemma-henselization@cref}{{[originallemma][0][]0A02}{[1][0][]source}{}{}{}}
\makeatother
\crefname{originallemma}{original lemma}{original lemmas}
\Crefname{originallemma}{Original lemma}{Original lemmas}
\makeatletter
\newlabel{sites-cohomology-lemma-cohomology-of-open}{{03F3}{source}{}{sourcepointer-2}{}}
\newlabel{sites-cohomology-lemma-cohomology-of-open@cref}{{[originallemma][0][]03F3}{[1][0][]source}{}{}{}}
\makeatother
\crefname{originalreference}{original reference}{original references}
\Crefname{originalreference}{Original reference}{Original references}
\makeatletter
\newlabel{theorem-colimit}{{09YQ}{source}{}{sourcepointer-3}{}}
\newlabel{theorem-colimit@cref}{{[originalreference][0][]09YQ}{[1][0][]source}{}{}{}}
\makeatother
\crefname{originalreference}{original reference}{original references}
\Crefname{originalreference}{Original reference}{Original references}
\makeatletter
\newlabel{theorem-gabber}{{09ZI}{source}{}{sourcepointer-4}{}}
\newlabel{theorem-gabber@cref}{{[originalreference][0][]09ZI}{[1][0][]source}{}{}{}}
\makeatother
\title{On a scheme with affine diagonal covered by n+1 affine opens the restriction of an injective torsion module has no cohomology above degree n}
\author{The Stacks Project Authors}
\date{stacks-project:etale-cohomology.tex:local-snapshot-20261011; GNU FDL 1.2 or later}
\begin{document}
\maketitle
\noindent\textit{Editorial note.} All delivered mathematical text is the author's own, copied from the pinned Stacks Project chapter file etale-cohomology.tex; no statement or proof step was rewritten, reordered or removed.\par\medskip
\noindent\textit{Editorial note.} The counted proof is the author's complete argument: the proof is written out in full: the affine case with the henselisation, Gabber's comparison and the colimit computation, and the induction step with the Mayer-Vietoris sequence and the affine-diagonal input.\par\medskip
\noindent\textit{Editorial note.} The label lemma-vanishing-restriction-injective names the counted statement, and the proof environment immediately following it in the source is delivered verbatim as its proof.\par\medskip
\noindent\textit{Editorial note.} The counted statement is carried by the author's own lemma environment in the chapter source.\par\medskip
\noindent\textit{Editorial note.} The statement and proof are copied from the pinned chapter file; the four cited inputs are kept as external pointers with their canonical Stacks tags.\par\medskip
\noindent\textit{Editorial note.} The preamble is the chapter's own native preamble excerpt already used for this source: the Xy-pic setup (native preamble lines 14-18) and the native theorem declarations with the mathematical macros (native preamble lines 154-217).\par\medskip
\noindent\textit{Allowed standard prerequisites (editorial).}\begin{itemize}
\item Henselisation of a pair and Gabber's comparison. More on Algebra, the description of the henselisation A\textasciicircum{}h as a filtered colimit of etale A-algebras with A/I = A\textasciicircum{}h/IA\textasciicircum{}h, together with the same chapter's henselian-pair comparison theorem. Supplies the identification H\textasciicircum{}q(Z, I|\_Z) = H\textasciicircum{}q(X\textasciicircum{}h, I|\_\{X\textasciicircum{}h\}) in the affine case.
\item Cohomology of filtered colimits and injectivity on an etale neighbourhood. Same chapter, the colimit theorem for etale cohomology and the lemma that the restriction of an injective sheaf along an etale morphism is injective. Supplies the colimit presentation of H\textasciicircum{}q(X\textasciicircum{}h, I|\_\{X\textasciicircum{}h\}) and the vanishing of each term of the colimit.
\item Mayer-Vietoris for an open cover with affine intersection. Same chapter, the Mayer-Vietoris exact sequence for the cover X = U union V used in the induction. Supplies the exact sequence that propagates the vanishing from U, V and U cap V to X.
\end{itemize}\medskip
\ifdefined\hypertarget\hypertarget{sourcepointer-1}{}\fi
\noindent\textit{Source reference.} Original lemma 0A02: sources/stacks-chow/tags.txt, source lines 12620--12620. Prerequisite pointer invoked at proof line(s) 15273; it names a standard fact of a sibling chapter that the counted argument uses, and every step of the counted proof built on it is retained verbatim in the delivered body.\par\medskip
\ifdefined\hypertarget\hypertarget{sourcepointer-2}{}\fi
\noindent\textit{Source reference.} Original lemma 03F3: sources/stacks-chow/tags.txt, source lines 4478--4478. Prerequisite pointer invoked at proof line(s) 15289; it names a standard fact of a sibling chapter that the counted argument uses, and every step of the counted proof built on it is retained verbatim in the delivered body.\par\medskip
\ifdefined\hypertarget\hypertarget{sourcepointer-3}{}\fi
\noindent\textit{Source reference.} Original reference 09YQ: sources/stacks-etale-cohomology/pinned-original-source.tex, source lines 7400--7400. Prerequisite pointer invoked at proof line(s) 15280; it names a standard fact of the same chapter that the counted argument uses, and every step of the counted proof built on it is retained verbatim in the delivered body.\par\medskip
\ifdefined\hypertarget\hypertarget{sourcepointer-4}{}\fi
\noindent\textit{Source reference.} Original reference 09ZI: sources/stacks-etale-cohomology/pinned-original-source.tex, source lines 15205--15205. Prerequisite pointer invoked at proof line(s) 15275; it names a standard fact of the same chapter that the counted argument uses, and every step of the counted proof built on it is retained verbatim in the delivered body.\par\medskip
\begin{lemma}
\label{lemma-vanishing-restriction-injective}
Let $X$ be a scheme with affine diagonal which can be covered by
$n + 1$ affine opens. Let $Z \subset X$ be a closed subscheme.
Let $\mathcal{A}$ be a torsion sheaf of rings on $X_\etale$
and let $\mathcal{I}$ be an injective sheaf of $\mathcal{A}$-modules
on $X_\etale$.
Then $H^q_\etale(Z, \mathcal{I}|_Z) = 0$ for $q > n$.
\end{lemma}

\begin{proof}
We will prove this by induction on $n$. If $n = 0$, then $X$ is affine.
Say $X = \Spec(A)$ and $Z = \Spec(A/I)$. Let $A^h$ be the filtered colimit
of \'etale $A$-algebras $B$ such that $A/I \to B/IB$ is an isomorphism.
Then $(A^h, IA^h)$ is a henselian pair and $A/I = A^h/IA^h$, see
More on Algebra, Lemma \ref{more-algebra-lemma-henselization}
and its proof. Set $X^h = \Spec(A^h)$.
By Theorem \ref{theorem-gabber}
we see that
$$
H^q_\etale(Z, \mathcal{I}|_Z) = H^q_\etale(X^h, \mathcal{I}|_{X^h})
$$
By Theorem \ref{theorem-colimit} we have
$$
H^q_\etale(X^h, \mathcal{I}|_{X^h}) =
\colim_{A \to B} H^q_\etale(\Spec(B), \mathcal{I}|_{\Spec(B)})
$$
where the colimit is over the $A$-algebras $B$ as above.
Since the morphisms $\Spec(B) \to \Spec(A)$ are \'etale,
the restriction $\mathcal{I}|_{\Spec(B)}$ is an injective
sheaf of $\mathcal{A}|_{\Spec(B)}$-modules
(Cohomology on Sites, Lemma \ref{sites-cohomology-lemma-cohomology-of-open}).
Thus the cohomology groups on the right are zero and we get the
result in this case.

\medskip\noindent
Induction step. We can use Mayer-Vietoris to do the induction step.
Namely, suppose that $X = U \cup V$ where $U$ is a union of $n$ affine
opens and $V$ is affine. Then, using that the diagonal of $X$ is affine,
we see that $U \cap V$ is the union of $n$ affine opens. Mayer-Vietoris
gives an exact sequence
$$
H^{q - 1}_\etale(U \cap V \cap Z, \mathcal{I}|_Z) \to
H^q_\etale(Z, \mathcal{I}|_Z) \to
H^q_\etale(U \cap Z, \mathcal{I}|_Z) \oplus
H^q_\etale(V \cap Z, \mathcal{I}|_Z)
$$
and by our induction hypothesis we obtain vanishing for $q > n$ as desired.
\end{proof}

\end{document}
